\chapter{Queries, intersection and overlap}
\label{ch:queries}

\section{Bounding boxes}

\fn{MFTLObjectAABB}\texttt{(obj, centre, half)} (collective) returns the axis-aligned box of the
owned vertices of an object as a centre and half-dimensions; \fn{MFTLPointInsideAABB} tests a
point against such a box with \texttt{MFTL\_VERYSMALL} slack. \mftl{} has no oriented bounding
boxes: the only consumer that needed a body box (\multi's immersed-boundary motion) uses an
axis-aligned one, and candidate searches use the LBVH below.

\section{Point in object}
\label{sec:pointinside}

\fn{MFTLPointInsideObject}\texttt{(obj, p)} classifies a point against a closed body by ray
casting (Figure~\ref{fig:raycast}): a ray from $p$ in a random direction (drawn on rank 0 and
broadcast) is intersected with the stored triangles of every rank, using the LBVH to find the
candidates; the number of crossings, summed over ranks, is odd for an inside point. If the ray
grazes a vertex or an edge anywhere, every rank retries with a new direction (at most 100
times; after that the point is taken as inside, as in \mftlone). The result is 1 (inside), 0
(outside) or $-1$ (on the surface). The call is collective and every rank passes the same
point. On a periodic domain the probe is re-imaged next to the body, which must be compact;
for a wrapping body use \fn{MFTLPhaseSetInside}.

\begin{figure}[ht]
\centering
\begin{tikzpicture}[scale=1]
  \draw[thick,fill=blue!6] plot[smooth cycle,tension=0.8] coordinates {(0,0) (2,-0.6) (3.6,0.2) (3.2,1.8) (1.6,1.4) (0.4,2.1) (-0.6,1.1)};
  \fill (1.2,0.8) circle (1.8pt) node[below left] {$p$};
  \draw[-{Stealth},red!70!black,thick] (1.2,0.8) -- (5.2,2.6) node[right] {ray};
  \fill (5.2,-0.2) circle (1.8pt) node[below] {$q$};
  \draw[-{Stealth},red!70!black,thick] (5.2,-0.2) -- (-1.2,0.5);

  \node[font=\small,align=left] at (7.3,1.1) {$p$: one crossing (inside)\\ $q$: two crossings (outside)};
\end{tikzpicture}
\caption{Point-in-object by ray casting: an odd number of crossings means inside.}
\label{fig:raycast}
\end{figure}

For many points against one body, a caller-held query builds the LBVH once:
\begin{lstlisting}[numbers=none]
MFTL_PointQuery *MFTLPointQueryCreate(MFTL_Object *obj);          /* collective */
int  MFTLPointQueryInside(MFTL_PointQuery *q, const double p[3]);
int  MFTLPointQueryInsideMany(MFTL_PointQuery *q, int n, const double *xyz, int *result);
void MFTLPointQueryDestroy(MFTL_PointQuery **q);
\end{lstlisting}
\fn{MFTLPointQueryInsideMany} classifies $n$ points (the same points on every rank) with the same
ray sequence and verdicts as $n$ single calls, but with $O(\text{retries})$ collectives instead of
$O(n)$. The query stays valid until the body changes. \textbf{Known defect:} on a non-periodic
domain the caller-held query deadlocks when some rank stores no triangle of the body; use it
only when every rank holds part of the body, as the example of \S\ref{sec:ex-queries} does
(\S\ref{sec:known-defects}).

\fn{MFTLPointTriangleDistance}\texttt{(obj, t, p)} returns the exact distance from a point to a
stored triangle (the \mftlone{} six-region classification); it uses the stored positions as
given, without minimum image, and is rank-local.

\section{The linear bounding-volume hierarchy}

\fn{MFTL\_GenerateLBVH}\texttt{(obj)} builds a bounding-volume hierarchy over all stored triangles
of an object (rank-local):
\begin{enumerate}
  \item the triangle centroids are normalised to their bounding box, quantised to $2^{10}$ cells
    per axis and interleaved into 30-bit Morton codes (Z-order);
  \item the (code, index) pairs are sorted with a radix sort;
  \item the tree is built top-down by splitting a range where the highest bit in which the keys
    of its first and last element differ changes; leaves hold at most four triangles;
  \item node boxes are refitted bottom-up, stored in single precision and inflated by
    $10^{-4}$.
\end{enumerate}
Nodes are 32 bytes (two per cache line). \fn{MFTL\_DestroyBVH} frees the tree.

\section{Triangle intersection}
\label{sec:intersection}

\fn{MFTL\_IntersectObjects}\texttt{(A, bvhA, B, bvhB, callback, user)} traverses two hierarchies
simultaneously and calls
\begin{lstlisting}[numbers=none]
typedef void (*IntersectionCallback)(MFTL_Object *objA, int triA, MFTL_Object *objB, int triB,
                                     void *user_data);
\end{lstlisting}
for every pair of stored triangles that intersect or come within $10^{-6}$ of each other. The
pair test is the Guigue--Devillers triangle test built on orientation predicates
(\fn{MFTL\_TriTriIntersect}), extended with nine edge--edge and six vertex--face distances for
the proximity. On a periodic domain the box tests shift by one period and each tested pair is
first placed in one common frame. Exactly coplanar pairs are reported as intersecting.

\fn{MFTLEnsembleIntersectObjects}\texttt{(ensA, ensB, callback, user)} is the collective driver:
it creates pairwise spatial ghosts, builds one LBVH per object and traverses every object pair
(pairs $i<j$ of one ensemble when \texttt{ensB} is \texttt{NULL}). The callback runs rank-locally
on stored triangles, so a pair near a rank boundary may be reported on two ranks.

\section{Overlap volumes}

\fn{MFTLEnsembleCalculateOverlap}\texttt{(ensA, ensB, \&records, \&n)} quantifies how much bodies
overlap, triangle pair by triangle pair (the \mftlone{} algorithm, with four documented
corrections). For each candidate pair the two triangles are projected on the plane normal to the
separation of the two bodies, the projected triangles are clipped against each other, and the
prism between the two sheets over the clip polygon gives a volume
$\tfrac13 A\sum(h_A-h_B)$ with exact zero crossings. Each record (\texttt{MFTL\_Overlap}) holds the
volume, the overlap areas on both sheets, the average surface normal
$\tfrac12(\vect{n}_A-\vect{n}_B)$, the centre of the overlap volume, an effective radius of
curvature and the global ids of the pair. The candidates are gathered with an all-gather
(owned triangles whose box meets the other body's box), every record is computed once, and the
full, deterministically ordered list is returned on every rank. On a periodic domain every body
must be compact and lie within half a period of its surface centre. \fn{MFTLCalculateOverlap}
computes a single pair. The 2-D helpers \fn{MFTLConvexHull2D}, \fn{MFTLPolygonArea2D} and
\fn{MFTLPolygonCentroid2D} are public as well.
